% Gesamttest «Polynomfunktionen» — Quelle des PDFs (python3 scripts/build-lp-pdf.py).
% Musterlösung und Raster: bewertungspaket.tex. Alle Werte mit python3 nachgerechnet (04.10.2026, nach der Prüfung überarbeitet).
% Kein \lt/\gt — pdfLaTeX kennt sie nicht, < und > tun es.
\documentclass[11pt]{article}
\usepackage{../lp-druck}
\renewcommand{\lpname}{Polynomfunktionen}
\renewcommand{\lpteilgebiet}{3.3}
\renewcommand{\lpbereich}{Schwerpunktfach}
\renewcommand{\lpfuss}{Gesamttest · 25 Punkte}
\begin{document}
\lpkopf{Gesamttest}{%
  \vspace{4pt}\noindent\begin{tabularx}{\linewidth}{@{}X X@{}}
  Code (kein Name): \hrulefill & Punkte: \hrulefill\ / 25
  \end{tabularx}}

\begin{lpinfo}{Anleitung}
\textbf{Zeit:} rund 30 Minuten. \textbf{Hilfsmittel:} Teil~A \textbf{ohne Taschenrechner} — die beiden Kompetenzen tragen im Lehrplan den Vermerk «auch ohne Hilfsmittel». Teil~B \textbf{mit Taschenrechner}. Lineal und Bleistift sind immer erlaubt.
Schreib zu jeder Aufgabe den \textbf{Rechenweg} auf — bewertet wird der Weg.
Danach: Lösung scannen oder fotografieren (ohne Namen und Standort) und mit dem \emph{Bewertungspaket} bewerten lassen.
\end{lpinfo}

\lpteil{Teil A · Nullstellen und Verlauf — ohne Taschenrechner}{Kapitel 1 bis 3 · 13 Punkte}

\begin{lpaufgabe}{G1}{Aus der Linearfaktordarstellung lesen}{4 P}
Gegeben ist $f(x) = -0.5\,(x+3)\,(x-2)^{3}$.
\begin{enumerate}[label=(\alph*),leftmargin=*,itemsep=1pt]
\item Gib Grad und Leitkoeffizient an.
\item Gib alle Nullstellen mit ihrer Vielfachheit an und beschreib, was der Graph an jeder Nullstelle tut.
\item Berechne $f(0)$.
\end{enumerate}
\lpplatz{7}
\end{lpaufgabe}

\begin{lpaufgabe}{G2}{Vom Graphen zur Gleichung}{4 P}
\begin{minipage}[t]{0.52\linewidth}
Der Graph gehört zu einer Polynomfunktion dritten Grades. Die markierten Punkte liegen auf Gitterpunkten.
Bestimme die Gleichung in Linearfaktordarstellung und begründe, welche Nullstelle doppelt ist.
\lpplatz{8}
\end{minipage}\hfill
\begin{minipage}[t]{0.42\linewidth}\vspace{0pt}\centering
\begin{tikzpicture}
\begin{axis}[width=6.4cm,height=7.6cm,axis lines=middle,xmin=-3.2,xmax=4.2,ymin=-4.5,ymax=8.5,
  xtick={-3,-2,-1,1,2,3,4},ytick={-4,-2,2,4,6,8},grid=both,grid style={lplinie},tick label style={font=\scriptsize},
  xlabel=$x$,ylabel=$y$,
  yticklabel style={fill=white,inner sep=1pt,xshift=-2pt},xticklabel style={fill=white,inner sep=1pt,yshift=-2pt}]
\addplot[lpblau,very thick,domain=-2.45:3.62,samples=160] {0.5*(x+2)*(x-2)^2};
\addplot[only marks,mark=*,color=lporange] coordinates {(-2,0) (2,0)};
\addplot[only marks,mark=*,color=black] coordinates {(0,4)};
\end{axis}
\end{tikzpicture}
\end{minipage}
\end{lpaufgabe}

\begin{lpaufgabe}{G3}{Den Verlauf beschreiben}{3 P}
Gegeben ist $g(x) = -x^{4} + 5x^{2} - 4$.
\begin{enumerate}[label=(\alph*),leftmargin=*,itemsep=1pt]
\item Wohin zeigen die beiden Enden des Graphen? Begründe mit Grad und Leitkoeffizient.
\item Welche Symmetrie hat der Graph? Begründe mit den Exponenten.
\item Wie viele Nullstellen und wie viele lokale Extremstellen kann $g$ höchstens haben?
\end{enumerate}
\lpplatz{6}
\end{lpaufgabe}

\begin{lpaufgabe}{G4}{Zweimal berühren}{2 P}
Der Graph einer Polynomfunktion hat genau zwei Nullstellen. An beiden berührt er die $x$-Achse, und beide Enden zeigen nach unten.
Welchen Grad hat die Funktion mindestens, und welches Vorzeichen hat ihr Leitkoeffizient? Begründe beides.
\lpplatz{5}
\end{lpaufgabe}

\lpteil{Teil B · Ausgezeichnete Stellen — mit Taschenrechner}{Kapitel 4 und 5 · 12 Punkte}

\begin{lpaufgabe}{G5}{Nullstellen berechnen}{4 P}
Berechne alle Nullstellen von $h(x) = x^{3} - 4x^{2} + x + 6$: Rate eine Nullstelle, spalte den Linearfaktor ab und löse den Quotienten. Gib $h$ in Linearfaktordarstellung an.
\lpplatz{10}
\end{lpaufgabe}

\begin{lpaufgabe}{G6}{Am Graphen ablesen}{3 P}
\begin{minipage}[t]{0.52\linewidth}
Abgebildet ist der Graph von $k(x) = -x^{3} + 3x^{2} + 1$. Hoch- und Tiefpunkt liegen auf Gitterpunkten.
\begin{enumerate}[label=(\alph*),leftmargin=*,itemsep=1pt]
\item Lies Hochpunkt und Tiefpunkt ab.
\item Bestimme das absolute Maximum und das absolute Minimum von $k$ auf $D = [-0.5;\, 3.5]$.
\end{enumerate}
\lpplatz{6}
\end{minipage}\hfill
\begin{minipage}[t]{0.42\linewidth}\vspace{0pt}\centering
\begin{tikzpicture}
\begin{axis}[width=6.4cm,height=7.6cm,axis lines=middle,xmin=-1.2,xmax=4.2,ymin=-6.5,ymax=7.5,
  xtick={-1,1,2,3,4},ytick={-6,-4,-2,2,4,6},grid=both,grid style={lplinie},tick label style={font=\scriptsize},
  xlabel=$x$,ylabel=$y$,
  yticklabel style={fill=white,inner sep=1pt,xshift=-2pt},xticklabel style={fill=white,inner sep=1pt,yshift=-2pt}]
\addplot[lpblau,very thick,domain=-1:3.55,samples=160] {-x^3+3*x^2+1};
\addplot[only marks,mark=*,color=lpgruen] coordinates {(0,1) (2,5)};
\end{axis}
\end{tikzpicture}
\end{minipage}
\end{lpaufgabe}

\begin{lpaufgabe}{G7}{Ein Extrempunkt exakt}{2 P}
Berechne den Extrempunkt von $p(x) = 0.5x^{2} - 3x + 2$ exakt und gib an, ob es ein Hoch- oder ein Tiefpunkt ist.
\lpplatz{5}
\end{lpaufgabe}

\begin{lpaufgabe}{G8}{Eine Flugbahn}{3 P}
Die Höhe eines Modellflugzeugs wird für $0 \le t \le 6$ beschrieben durch $h(t) = -t^{3} + 6t^{2}$ ($t$ in s, $h$ in m).
\begin{enumerate}[label=(\alph*),leftmargin=*,itemsep=1pt]
\item Berechne die Nullstellen von $h$ durch Ausklammern und gib ihre Vielfachheit an.
\item Bestimme mit einer Wertetabelle (Schrittweite $1$), wann das Flugzeug in der Tabelle am höchsten ist und wie hoch es dann fliegt.
\end{enumerate}
\lpplatz{8}
\end{lpaufgabe}
\end{document}
